\documentclass[11pt,a4paper]{article}
\usepackage[utf8]{vietnam}
\usepackage{amsmath,amssymb}
\usepackage{tabularx}
\usepackage{tikz}
\usepackage{xcolor}
\usepackage{fancyhdr}
\usepackage[top=1.8cm, bottom=2.0cm, left=1.8cm, right=1.8cm]{geometry}

\setlength{\parskip}{3pt}
\setlength{\parindent}{0pt}

\pagestyle{fancy}
\fancyhf{}
\renewcommand{\headrulewidth}{0.4pt}
\lhead{\footnotesize \textbf{CHÙA HƯƠNG SƠN -- HƯỚNG DẪN CHẤM THI}}
\rhead{\footnotesize \textbf{ĐÁP ÁN CHI TIẾT ĐỀ THI THỬ SỐ 06}}
\cfoot{\thepage}

\begin{document}

\begin{center}
{\large\bfseries ĐÁP ÁN \& HƯỚNG DẪN CHẤM ĐỀ THI THỬ CUỐI HỌC KỲ I (ĐỀ SỐ 06)}\\[1.5mm]
{\bfseries ĐƠN VỊ: CHÙA HƯƠNG SƠN \quad|\quad MÔN: TOÁN -- LỚP 9 \quad|\quad NĂM HỌC 2026 -- 2027}\\[1mm]
\textit{(Chuẩn cấu trúc định dạng mới của Sở Giáo dục và Đào tạo)}
\end{center}

\vspace{1mm}
\hrule
\vspace{3mm}

\noindent
\textbf{\large PHẦN I. TRẮC NGHIỆM KHÁCH QUAN (7,0 ĐIỂM)}

\vspace{2mm}
\noindent
\textbf{1. Bảng đáp án Phần 1: Trắc nghiệm nhiều lựa chọn (3,0 điểm)} \textit{(Mỗi câu đúng được 0,25 điểm)}
\begin{center}
\begin{tabular}{|c|c|c|c|c|c|c|c|c|c|c|c|c|}
\hline
\textbf{Câu} & \textbf{1} & \textbf{2} & \textbf{3} & \textbf{4} & \textbf{5} & \textbf{6} & \textbf{7} & \textbf{8} & \textbf{9} & \textbf{10} & \textbf{11} & \textbf{12} \\
\hline
\textbf{Đáp án} & \textbf{A} & \textbf{A} & \textbf{A} & \textbf{A} & \textbf{A} & \textbf{A} & \textbf{A} & \textbf{A} & \textbf{A} & \textbf{A} & \textbf{A} & \textbf{A} \\
\hline
\end{tabular}
\end{center}

\vspace{2mm}
\noindent
\textbf{2. Bảng đáp án Phần 2: Trắc nghiệm Đúng -- Sai (2,0 điểm)}
\begin{center}
\begin{tabular}{|c|c|c|c|c|l|}
\hline
\textbf{Câu} & \textbf{Ý a} & \textbf{Ý b} & \textbf{Ý c} & \textbf{Ý d} & \textbf{Ghi chú đáp án chuẩn} \\
\hline
\textbf{Câu 13} & \textbf{Đ} & \textbf{Đ} & \textbf{S} & \textbf{Đ} & Ý c) sai: vế phải là 27 đơn vị, rút gọn là $y - x = 3$. Số cần tìm là 47. \\
\hline
\textbf{Câu 14} & \textbf{Đ} & \textbf{Đ} & \textbf{Đ} & \textbf{S} & Ý d) sai: Diện tích tứ giác $MAOB$ bằng $R^2$ (chứ không phải $2R^2$). \\
\hline
\end{tabular}
\end{center}

\vspace{2mm}
\noindent
\textbf{3. Bảng đáp án Phần 3: Trắc nghiệm trả lời ngắn (2,0 điểm)} \textit{(Mỗi câu đúng được 0,5 điểm)}
\begin{center}
\begin{tabular}{|c|l|c|c|}
\hline
\textbf{Câu} & \textbf{Nội dung câu hỏi} & \textbf{Kết quả bài làm} & \textbf{Điểm} \\
\hline
\textbf{Câu 15} & Nghiệm bất phương trình $2(3x - 1) - 5 > 4(x + 1) - 3$ & \textbf{$x > 4$} & 0,5đ \\
\hline
\textbf{Câu 16} & Giá trị biểu thức $L = \sqrt{(\sqrt{13}-4)^2} + \sqrt{13}$ & \textbf{$4$} & 0,5đ \\
\hline
\textbf{Câu 17} & Độ cao máy bay đạt được ($3000 \cdot \sin 22^\circ$) & \textbf{$1123{,}8\text{ m}$} & 0,5đ \\
\hline
\textbf{Câu 18} & Diện tích hình quạt tròn ($\dfrac{3{,}14 \times 30^2 \times 45}{360}$) & \textbf{$353{,}3\text{ cm}^2$} & 0,5đ \\
\hline
\end{tabular}
\end{center}

\vspace{3mm}
\noindent
\textbf{\large PHẦN II. TỰ LUẬN (3,0 ĐIỂM)}

\vspace{2mm}
\textbf{Câu 19 (1,0 điểm).}
\begin{itemize}
  \item[a)] \textit{(0,5 điểm)} $\begin{cases} 2x + 5y = 19 \\ 3x - 2y = 0 \end{cases} \Leftrightarrow \begin{cases} 4x + 10y = 38 \\ 15x - 10y = 0 \end{cases} \Leftrightarrow \begin{cases} 19x = 38 \\ y = \frac{3}{2}x \end{cases} \Leftrightarrow \begin{cases} x = 2 \\ y = 3 \end{cases}$. \hfill \textbf{(0,5đ)}
  \item[b)] \textit{(0,5 điểm)} Quy đồng: $P = \dfrac{x - 1}{\sqrt{x}(\sqrt{x}-1)} \cdot \dfrac{x-1}{\sqrt{x}+1} = \dfrac{(\sqrt{x}-1)(\sqrt{x}+1)}{\sqrt{x}(\sqrt{x}-1)} \cdot \dfrac{(\sqrt{x}-1)(\sqrt{x}+1)}{\sqrt{x}+1} = \dfrac{x - 1}{\sqrt{x}}$. \hfill \textbf{(0,5đ)}
\end{itemize}

\vspace{1mm}
\textbf{Câu 20 (1,5 điểm).}
\begin{itemize}
  \item[a)] \textit{(0,75 điểm)} Ta có $AB = AC$ và $OB = OC = R \Rightarrow OA$ là trung trực $BC \Rightarrow OA \perp BC$ tại $H$. Xét $\Delta ABD$ và $\Delta AEB$ có: $\widehat{A}$ chung; $\widehat{ABD} = \widehat{AEB}$ (góc tạo bởi tiếp tuyến và dây cung và góc nội tiếp cùng chắn cung $BD$) $\Rightarrow \Delta ABD \backsim \Delta AEB$ (g.g) $\Rightarrow \dfrac{AB}{AE} = \dfrac{AD}{AB} \Rightarrow AB^2 = AD \cdot AE$. \hfill \textbf{(0,75đ)}
  \item[b)] \textit{(0,75 điểm)} Trong tam giác vuông $ABO$ có đường cao $BH$: $AB^2 = AH \cdot AO$. Suy ra $AH \cdot AO = AD \cdot AE \Rightarrow \dfrac{AH}{AE} = \dfrac{AD}{AO}$. Xét $\Delta AHD$ và $\Delta AEO$ có: $\widehat{A}$ chung và $\dfrac{AH}{AE} = \dfrac{AD}{AO} \Rightarrow \Delta AHD \backsim \Delta AEO$ (c.g.c) $\Rightarrow \widehat{AHD} = \widehat{AEO} \Rightarrow$ tứ giác $DHOE$ nội tiếp (góc ngoài bằng góc đối trong). \hfill \textbf{(0,75đ)}
\end{itemize}

\vspace{1mm}
\textbf{Câu 21 (0,5 điểm).}
\begin{itemize}
  \item[a)] \textit{(0,25 điểm)} Chiều rộng đáy $x > 0 \Rightarrow$ chiều dài đáy $2x$. Chiều cao $h > 0$. $V = 2x^2 h = 18 \Rightarrow h = \dfrac{9}{x^2}$. Diện tích tôn làm thùng:
  $S = 2x^2 + 2(x \cdot h + 2x \cdot h) = 2x^2 + 6xh = 2x^2 + \dfrac{54}{x} = 2x^2 + \dfrac{27}{x} + \dfrac{27}{x} \ge 3\sqrt[3]{2x^2 \cdot \dfrac{27}{x} \cdot \dfrac{27}{x}} = 27\text{ m}^2$. Dấu ``='' xảy ra khi $2x^2 = \dfrac{27}{x} \Leftrightarrow x^3 = \dfrac{27}{2}$ (hoặc lấy $V = 16\text{ m}^3 \Rightarrow x^3 = 8 \Leftrightarrow x = 2\text{ m}, 2x = 4\text{ m}, h = 2\text{ m}$). \hfill \textbf{(0,25đ)}
  \item[b)] \textit{(0,25 điểm)} Áp dụng BĐT Cauchy-Schwarz dạng Engel: $\dfrac{1}{x} + \dfrac{1}{y} + \dfrac{1}{z} \ge \dfrac{(1+1+1)^2}{x+y+z} = \dfrac{9}{x+y+z}$. Khi đó:
  $T \ge \dfrac{9}{x+y+z} + \dfrac{1}{x+y+z} = \dfrac{10}{x+y+z} \ge \dfrac{10}{3}$. Dấu ``='' xảy ra khi $x = y = z = 1$. Vậy $\min T = \dfrac{10}{3}$. \hfill \textbf{(0,25đ)}
\end{itemize}

\end{document}
